\documentclass[10pt,a4paper]{amsart}
\usepackage[margin=19mm]{geometry}
\usepackage{amsmath,amssymb,mathtools}
\usepackage[scr=rsfs,cal=euler]{mathalfa}
\usepackage{xfrac}
\usepackage[unicode,hidelinks,bookmarks=false]{hyperref}
\newcommand{\ie}{i.e.\ }
\newcommand{\cf}{cf.\ }
\newcommand{\resp}{respectively\ }
\newcommand{\Subsection}{\subsection}
\providecommand{\nobreakdash}{\nobreak\hbox{-}}
\setlength{\emergencystretch}{2em}
\multlinegap0pt
\usepackage{amssymb}
\usepackage{enumerate}
\usepackage{algorithm}\usepackage{algpseudocode}
\usepackage{bm}
\usepackage{xcolor}
\usepackage{cancel}
\newtheorem{prop}{Proposition}
\title{Conjugate phase retrieval in shift-invariant spaces generated by a Gaussian}
\author{Yang Chen, Cheng Cheng}
\date{Source: arXiv:2412.03807v2 (CC BY 4.0)}
\begin{document}
\maketitle

\noindent\emph{Source and licence.} Conjugate phase retrieval in shift-invariant spaces generated by a Gaussian. Version arXiv:2412.03807v2 (CC BY 4.0), distributed by arXiv under the Creative Commons Attribution 4.0 International licence, http://creativecommons.org/licenses/by/4.0/, as recorded in the arXiv licence metadata for this exact version and shown on the abstract page https://arxiv.org/abs/2412.03807v2. Full licence notice: \texttt{licenses/arxiv-2412.03807-CC-BY-4.0.txt}.\par\smallskip

\noindent\textit{Editorial scope.} Counted unit: the article's prop hermite.lem (source0.tex lines 686--688). The article's complete original proof of it is delivered with it (source0.tex lines 690--740, one \texttt{proof} environment of 51 lines). No mathematics is written, repaired or re-derived: the statement and the proof are the article's own lines, in the article's own notation, and no retained line is substituted or re-flowed.  Retained uncounted as the source-local context the proof needs: lines 647--654; lines 670--676.  Nothing was omitted from any retained span, so the package carries no omission record.  No retained line contains a citation, so the wrapper carries no bibliography and no bibliography entry.  No retained line contains a figure environment, an image-inclusion command or drawing code, so no image is delivered and the package declares no asset dependency.  Editorial formatting only, in addition: an AMS \texttt{amsart} preamble that restates the article's own declarations and macros used by the retained lines, namely the environments \texttt{prop} on separate counters, together with the standard packages those lines need.  The retained environments print in this document as Proposition 1 in order of appearance, whereas the article prints the same items with the numbering of its own full document.  The complete native source of the article as distributed in the arXiv e-print for this exact version is bundled unchanged as \texttt{sources/169636/source0.tex}.
\begin{equation}\label{h.sis}
V(h_n)
:=
\left\{
\sum_{k\in\mathbb Z} c_k h_n(\cdot-k): \;
\{c_k\}_{k\in\mathbb Z}\in\ell^\infty(\mathbb Z)
\right\}.
\end{equation}

\begin{equation}\label{hermite.sis.expansion}
f(x)
=
\zeta_n\sum_{l=0}^{n}(-1)^l(\sqrt2)^l m_l
\sum_{k\in\mathbb Z} c_k
\frac{{\rm d}^l}{{\rm d}x^l}\big(e^{-(x-k)^2/2}\big).
\end{equation}

\begin{prop}\label{hermite.lem}
{\rm Let $f\in V(h_n)$ be a function in the shift-invariant space generated by the Hermite function $h_n$ of order $n$ in \eqref{h.sis}.  Then the function $f$  and its squared modulus function $|f|^2$ can be extended to  the entire function of order at most 2.  }
\end{prop}

\begin{proof} 
Let $f\in V(h_n)$. By the expansion \eqref{hermite.sis.expansion}, the function $f$ can be written
as a finite linear combination of derivatives of a Gaussian shift-invariant signal. 
Set
\[
F(z):=\sum_{k\in\mathbb Z} c_k e^{-(z-k)^2/2}, \qquad z=x+iy\in\mathbb C,
\]
where $\{c_k\}_{k\in\mathbb Z}\in\ell^\infty(\mathbb Z)$.

For each $k\in\mathbb Z$,
\[
|e^{-(z-k)^2/2}|
= e^{(y^2-(x-k)^2)/2}
\le e^{y^2/2} e^{-(x-k)^2/2}.
\]
Hence, for any strip $\{z\in\mathbb C:|y|\le R\}$,
\[
|F(z)|
\le \|c\|_\infty e^{R^2/2}
\sum_{k\in\mathbb Z} e^{-(x-k)^2/2}.
\]
Since the sum $\sum_{k\in\mathbb Z} e^{-(x-k)^2/2}$ is uniformly bounded in
$x\in\mathbb R$, the series defining $F$ converges absolutely and uniformly on compact subsets of $\mathbb C$. Consequently, $F$ is entire and satisfies
\[
|F(z)|\le C e^{|z|^2},
\]
so that $F$ is an entire function of order at most $2$.

Moreover, it can be infinitely differentiable. For any $l\ge 0$,
\[
F^{(l)}(z)=\sum_{k\in\mathbb Z} c_k
\frac{{\rm d}^l}{{\rm d}z^l}\big(e^{-(z-k)^2/2}\big),
\]
and  each $F^{(l)}$ is also an entire function of order at most $2$.
Combining with \eqref{hermite.sis.expansion}, we conclude that $f$ admits the analytic extension
\begin{equation}\label{f.rep}
f(z)=\zeta_n\sum_{l=0}^n (-1)^l(\sqrt2)^l m_l\,F^{(l)}(z),
\end{equation}
which is an entire function of order no more than $2$. 

\smallskip 

Define $G(z):=f(z)\,\overline{f(\overline{z})}$. Since both
$f(z)$ and $z\mapsto\overline{f(\overline{z})}$ are entire, $G$ is entire. On the real axis,
\[
G(x)=f(x)\overline{f(x)}=|f(x)|^2.
\]
Moreover, $G$ has order at most $2$ because it is a product of two entire
functions of order at most $2$. Hence $|f|^2$ admits an entire extension of order at most $2$.  This completes the proof.

\end{proof}

\end{document}
